% ===================================================================
%  TABLEAU N° 7 — Le village en hiver. La maison de campagne au
%                 printemps.
%
%  Traduction, faite en 2026, en francais standard du Quebec, sur
%  l'ido de texto/io. Regles de la colonne : voir 10-tableau-01.tex.
%  Deux scenes, deux numerotations : les cles portent c1 et c2.
%
%  LES SIX CHOIX PROPRES A CE TABLEAU :
%    (37) tas de neige ............ banc de neige
%    (35) cache-nez ............... foulard
%    (13) institutrice ............ enseignante
%    (15) instituteur ............. enseignant
%    (32) servante ................ domestique
%    (79) niche (a lapins) ........ clapiers
%
%  « BANC DE NEIGE » EST LE MOT DU PAYS, ET IL EST PLUS PRECIS QUE
%  CELUI QU'IL REMPLACE. Le fac-simile ecrit « un tas de neige » et
%  l'ido « niv-amaso », un amas ; mais l'amas de la gravure est celui
%  que le vent et la pelle ont pousse contre le pont, et c'est
%  exactement ce qu'un banc de neige designe ici. Le bonhomme de
%  neige, lui, ne change pas : le Quebec le nomme comme la France.
%
%  CINQ MOTS DU FAC-SIMILE RESTENT PARCE QUE LE QUEBEC LES DIT
%  ENCORE : barbier (les salons de barbier n'ont pas disparu ici),
%  glissoire, traineau, bonhomme de neige, marrons.
%
%  LE REPAS DE L'ALINEA 4 EST UN DINER. L'ido dit « dejuno » et le
%  fac-simile « dejeuner » ; c'est le repas du milieu du jour, et au
%  Quebec c'est le diner. Regle posee au tableau 4.
%
%  « SERVANTE » PART, ET C'EST LA TROISIEME FOIS. Comme la « bonne »
%  du tableau 4 et le « domestique » du tableau 6, la personne au
%  service d'une maison se dit ici domestique : l'ido n'a qu'un mot
%  (« servistino ») et le francais de 1926 en avait trois, tous
%  marques. On garde le seul qui ne le soit plus.
%
%  LA PLAISANTERIE DU CORDONNIER TIENT A DEUX MOTS, ET IL FALLAIT LES
%  DEUX. L'ido oppose « shuifisto », qui FAIT la chaussure, a
%  « shureparisto », qui la REPARE — en ville le premier n'avait pas
%  de clients, au village le second ne chome pas. Rendre les deux par
%  « cordonnier » aurait efface la chute. D'ou bottier et cordonnier.
% ===================================================================

%%K t07-noto-1 noto
\VUnotes{143.2mm}{%
(*) Voir, pour le détail des repas, les tableaux 6 et 13.%
}

%%K t07-apar-1 apar
\VUsaut{4.0mm}
\VUcentre{13.2pt}{90}{DEUXIÈME SÉRIE}
\VUsaut{3.0mm}
\VUfilet{16mm}
\VUsaut{5.6mm}
\VUtitre{15.0pt}{60}{TABLEAU N\textsuperscript{o} 7}
\VUsaut{3.0mm}

%%K t07-tit-1 sub
\VUcentre{12.0pt}{0}{\textit{Première scène.}}
\VUsaut{3.0mm}

%%K t07-tit-2 sub
\VUpk{10.2pt}{40}{Le village en hiver.}
\VUsaut{4.0mm}

%%K t07-c1-01-1 p
1. --- Pendant les vacances du jour de l'An, j'ai accompagné mon père à
Villeneuve, un petit village où il avait quelques affaires de commerce à régler.

%%K t07-c1-02-1 p
2. --- Nous avons pris la \VUgras{diligence} \textsuperscript{(11)} qui fait le
service entre la ville et ce bourg; mais comme il faisait très froid, ce voyage
dans une voiture peu confortable n'a pas été très agréable.

%%K t07-c1-03-1 p
3. --- Aussi avons-nous éprouvé un vrai plaisir quand nous avons vu le
\VUgras{conducteur} arrêter sa voiture sur la place, devant l'\VUgras{auberge}
\textsuperscript{(2)} de \textit{l'Ours blanc}. L'\VUgras{aubergiste}
\textsuperscript{(4)} nous attendait sur le seuil de sa maison, en fumant
tranquillement sa \VUgras{pipe}.

%%K t07-c1-03-2 p
Il nous a invités à entrer et je ne me suis pas fait prier pour m'installer
devant le feu de sarments qui pétillait dans l'âtre. Je me suis alors bien moqué
du \VUgras{vent du nord} qui faisait grincer la vieille \VUgras{enseigne}
\textsuperscript{(3)} et emportait en noirs tourbillons \VUgras{la fumée}
\textsuperscript{(1)} qui sortait de la cheminée.

%%K t07-c1-04-1 p
4. --- C'était l'heure du dîner. Sans attendre plus longtemps, nous avons fait
honneur au repas simple mais savoureux qu'on nous a servi (*).

%%K t07-tit-3 sub
\VUpk{10.2pt}{40}{Quelques visites.}
\VUsaut{3.0mm}

%%K t07-c1-05-1 p
5. --- Après le dîner, j'ai accompagné mon père, qui est assureur, chez ses
clients. Nous avons d'abord visité le \VUgras{maréchal-ferrant}
\textsuperscript{(5)}, qui habite près de l'auberge. Il était occupé à ferrer un
cheval. J'ai admiré la force de son compagnon, qui maintenait solidement sur son
tablier de cuir le \VUgras{sabot} du cheval, pendant que le maréchal fixait le
\VUgras{fer} \textsuperscript{(6)} à grands coups de marteau.

%%K t07-c1-06-1 p
6. --- Nous sommes ensuite entrés chez le \VUgras{cordonnier}
\textsuperscript{(7)} du village, le vieux Jacques. Il a beau avoir pour
enseigne une magnifique \VUgras{botte}, il se contentait de ressemeler une
vieille paire de souliers percés.

%%K t07-c1-06-2 p
Le vieux nous a dit en riant : « En ville, j'étais "bottier" et je n'avais pas
de clients. Ici, je suis "cordonnier" et l'ouvrage ne manque pas. »

%%K t07-c1-07-1 p
7. --- Il nous a parlé de son voisin le \VUgras{barbier} \textsuperscript{(8)},
dont la clientèle n'est pas contente. Ses \VUgras{rasoirs} sont toujours
ébréchés et ses \VUgras{ciseaux} ne coupent jamais bien. Mais comme il est le
seul barbier du village, on est bien obligé de se faire raser et couper les
cheveux par lui. C'est d'ailleurs un homme avare et désagréable.

%%K t07-c1-08-1 p
8. --- Heureusement, les autres \VUgras{voisins} ne lui ressemblent pas. Ainsi
le \VUgras{charron} \textsuperscript{(9)} est un homme bon, travaillant et
serviable. C'est un plaisir de lui faire réparer une voiture. Son travail est
toujours bien fini.

%%K t07-c1-09-1 p
9. --- Le vieux Jacques ne se plaignait pas non plus du \VUgras{sellier}
\textsuperscript{(10)}, qu'il aide parfois à coudre les \VUgras{harnais} quand
le travail presse.

%%K t07-c1-09-2 p
Mon père lui a demandé des nouvelles de sa famille : « Tout le monde va bien. Ma
petite-fille est à l'\VUgras{école} \textsuperscript{(12)}. Son
\VUgras{enseignante} \textsuperscript{(13)} est très contente d'elle, c'est une
bonne \VUgras{élève} \textsuperscript{(14)}. Elle ne ressemble pas à mon
garnement de fils; lui, il ne pense qu'à jouer. L'\VUgras{enseignant}
\textsuperscript{(15)} n'arrive à rien lui faire apprendre. »

%%K t07-tit-4 sub
\VUsaut{3.0mm}
\VUpk{10.2pt}{40}{Bavardages de village.}
\VUsaut{3.0mm}

%%K t07-c1-10-1 p
10. --- Le vieux nous a ensuite raconté les nouvelles du village. On va
construire une nouvelle école, parce que celle qui est installée dans la
\VUgras{mairie} est devenue trop petite. Ce sera d'ailleurs facile, parce qu'il
suffira d'acheter une partie du jardin du \VUgras{meunier}. Ce sera une bonne
affaire pour le pauvre homme. À cause du froid, les \VUgras{meules} du
\VUgras{moulin} \textsuperscript{(16)} ne peuvent plus moudre le blé, parce que
la \VUgras{roue} \textsuperscript{(17)} a été bloquée par la \VUgras{glace}
\textsuperscript{(25)} du \VUgras{ruisseau} \textsuperscript{(23)}.

%%K t07-c1-11-1 p
11. --- « À propos de construction, avez-vous vu notre nouvelle
\VUgras{église} \textsuperscript{(18)}? Elle est bien pittoresque sous son
manteau de neige. Les \VUgras{corbeaux} \textsuperscript{(19)}, qui ne
reconnaissent plus leur vieux clocher, se perchent tristement sur le grand arbre
du \VUgras{cimetière} \textsuperscript{(20)}. »

%%K t07-c1-12-1 p
12. --- Cette idée du cimetière a rappelé au cordonnier les personnes que mon
père avait connues et qui sont mortes l'an dernier. « Que de \VUgras{tombes}
\textsuperscript{(21)}, mon bon monsieur, se sont ajoutées aux autres, sous les
\VUgras{cyprès} \textsuperscript{(22)}! La mort a fauché notre vieux notaire.
Tous les villageois ont assisté à ses \VUgras{funérailles}. »

%%K t07-c1-13-1 p
13. --- « Voilà son successeur qui patine sur le ruisseau. Il vient justement de
me faire réparer la courroie de son \VUgras{patin} \textsuperscript{(26)}, qui
était cassée. »

%%K t07-c1-14-1 p
14. --- « Voilà aussi, sur la rive du ruisseau, le nouveau \VUgras{garde
champêtre} \textsuperscript{(27)}. C'est un ancien gendarme qui terrorise les
garnements du village. Ils se sauvent dès qu'ils aperçoivent ses
\VUgras{guêtres} \textsuperscript{(29)} et son \VUgras{képi}
\textsuperscript{(28)}. »

%%K t07-c1-15-1 p
15. --- « Ah! voilà le vieux Joseph qui amène une \VUgras{charrette de fagots}
\textsuperscript{(30)} pour le boulanger. --- C'est vrai, a dit mon père, je
reconnais son attelage de \VUgras{mules} \textsuperscript{(31)}. J'ai justement
besoin de lui parler, je vais en profiter. Au revoir, père Jacques. »

%%K t07-c1-16-1 p
16. --- Juste au moment où nous sortions, les enfants du village quittaient
l'école et s'élançaient en courant vers la \VUgras{glissoire}
\textsuperscript{(33)}, au risque de tomber et de se casser un membre dans leur
\VUgras{chute} \textsuperscript{(34)}. L'un d'eux a pris son \VUgras{traîneau}
\textsuperscript{(32)} chez le charron, y a assis un de ses camarades et est
parti en courant.

%%K t07-c1-17-1 p
17. --- D'autres se sont précipités vers un \VUgras{banc de neige}
\textsuperscript{(37)}, près du \VUgras{pont} \textsuperscript{(40)}, et se sont
mis à bombarder le \VUgras{bonhomme de neige} \textsuperscript{(39)} qu'ils
avaient élevé la veille et à qui ils avaient mis un chapeau, une pipe et un sac
d'école en bandoulière.

%%K t07-c1-18-1 p
18. --- Le vent glacial et le froid les inquiètent bien peu. La plupart n'ont
même pas de \VUgras{foulard} \textsuperscript{(35)} et, pour se réchauffer après
la bataille de \VUgras{boules de neige} \textsuperscript{(38)}, il leur suffit
d'une poignée de \VUgras{marrons} \textsuperscript{(42)} bien chauds achetés à
la vieille \VUgras{marchande} Jacqueline, qui grelotte de froid sous son
\VUgras{châle} \textsuperscript{(43)} trop mince.

%%K t07-tit-5 sub
\VUsaut{3.0mm}
\VUcentre{12.0pt}{0}{\textit{Deuxième scène.}}
\VUsaut{3.0mm}

%%K t07-tit-6 sub
\VUpk{10.2pt}{40}{La maison de campagne au printemps.}
\VUsaut{4.0mm}

%%K t07-c2-01-1 p
1. --- Malgré tous les plaisirs de l'hiver, c'est avec une joie profonde que
nous saluons le retour du \VUgras{printemps}.

%%K t07-c2-01-2 p
Quel tableau réjouissant que celui d'une cour de ferme, le soir, au mois de mai!

%%K t07-c2-02-1 p
2. --- À l'horizon, le \VUgras{soleil} \textsuperscript{(1)} se couche
lentement, inondant la \VUgras{campagne} \textsuperscript{(3)} de ses
\VUgras{rayons} \textsuperscript{(2)} de pourpre. Le \VUgras{laboureur}
\textsuperscript{(6)} diligent se hâte de labourer son \VUgras{champ}
\textsuperscript{(4)}.

%%K t07-c2-02-2 p
Avec sa \VUgras{charrue} \textsuperscript{(5)} attelée d'un vigoureux
\VUgras{cheval} \textsuperscript{(7)}, il trace le \VUgras{sillon}
\textsuperscript{(8)} dans lequel le \VUgras{semeur} \textsuperscript{(9)}
jettera à pleines mains la \VUgras{semence} qui produira les épis dorés.

%%K t07-c2-03-1 p
3. --- Les \VUgras{faucheurs} \textsuperscript{(11)}, munis de leurs faux, ont
fauché l'herbe du pré, les faneurs ont fait sécher le \VUgras{foin}
\textsuperscript{(10)} en le retournant avec les \VUgras{fourches}
\textsuperscript{(12)} et les râteaux, et les voilà qui reviennent.

%%K t07-c2-03-2 p
Les uns marchent devant la lourde voiture tirée par des bœufs; ils portent leur
outil sur l'épaule. D'autres, plus paresseux, se sont confortablement installés
sur le foin odorant.

%%K t07-c2-04-1 p
4. --- En passant, ils font un salut amical au \VUgras{jardinier}
\textsuperscript{(16)}, occupé dans le \VUgras{verger} \textsuperscript{(13)} à
tailler les \VUgras{arbres fruitiers} et à greffer les \VUgras{poiriers}
\textsuperscript{(14)}. Les \VUgras{pruniers} \textsuperscript{(15)} sont en
fleurs et, dans le \VUgras{potager} \textsuperscript{(17)} bien fumé, les
légumes, les choux et les salades sont déjà en belle santé.

%%K t07-c2-04-2 p
Sur le petit mur, un beau \VUgras{paon} \textsuperscript{(18)} déploie sa queue
pour faire admirer ses magnifiques plumes.

%%K t07-tit-7 sub
\VUpk{10.2pt}{40}{La cour de ferme.}
\VUsaut{3.0mm}

%%K t07-c2-05-1 p
5. --- Les portes de l'\VUgras{écurie} \textsuperscript{(19)} sont ouvertes et
l'on aperçoit le râtelier et la \VUgras{litière} \textsuperscript{(23)} de la
\VUgras{jument} \textsuperscript{(21)} que le \VUgras{palefrenier}
\textsuperscript{(20)} vient de dételer. Le \VUgras{poulain}
\textsuperscript{(22)}, si comique avec ses longues jambes, suit sa mère en
gambadant.

%%K t07-c2-06-1 p
6. --- Près du \VUgras{puits} \textsuperscript{(29)}, les \VUgras{vaches}
\textsuperscript{(25)} vont boire dans l'\VUgras{auge} \textsuperscript{(26)} de
bois qui leur sert d'abreuvoir. Le \VUgras{veau} \textsuperscript{(24)} en
profite pour téter, et le \VUgras{bœuf} \textsuperscript{(27)}, qui sent la
bonne odeur du fourrage, mugit joyeusement et dresse fièrement sa tête aux
\VUgras{cornes} \textsuperscript{(28)} puissantes.

%%K t07-c2-07-1 p
7. --- Tout le monde travaille. La \VUgras{domestique} \textsuperscript{(32)}
tient en main la \VUgras{corde} \textsuperscript{(31)} du puits. Elle puise de
l'eau dans un \VUgras{seau} \textsuperscript{(30)} pour faire boire les bêtes.
Auprès de la \VUgras{grange} \textsuperscript{(33)}, un homme ratisse la paille,
et devant la porte de la \VUgras{maison} \textsuperscript{(34)} une vieille
\VUgras{fileuse} \textsuperscript{(35)}, avec son rouet et sa
\VUgras{quenouille}, file le \VUgras{lin} ou le \VUgras{chanvre} comme
autrefois.

%%K t07-tit-8 sub
\VUsaut{3.0mm}
\VUpk{10.2pt}{40}{Le fermier.}
\VUsaut{3.0mm}

%%K t07-c2-08-1 p
8. --- Le \VUgras{fermier} \textsuperscript{(36)} dirige tous ces travaux et
donne ses instructions à ses employés. Il recommande de rentrer la
\VUgras{herse} \textsuperscript{(40)} et la \VUgras{pioche}
\textsuperscript{(39)} qu'on a oubliées près de la \VUgras{niche}
\textsuperscript{(38)} du \VUgras{chien} \textsuperscript{(37)}.

%%K t07-c2-09-1 p
9. --- Ses cris effarouchent un \VUgras{pigeon} \textsuperscript{(41)}, qui
s'enfuit à tire-d'aile vers le \VUgras{pigeonnier} \textsuperscript{(42)}, et
les \VUgras{poules} \textsuperscript{(43)}, qui se hâtent de rentrer au
\VUgras{poulailler} \textsuperscript{(44)}.

%%K t07-c2-10-1 p
10. --- Devant la \VUgras{bergerie} \textsuperscript{(48)}, voici le vieux
\VUgras{berger} \textsuperscript{(45)} qui ramène son \VUgras{troupeau}
\textsuperscript{(49)}. Tenant sa \VUgras{houlette} \textsuperscript{(46)},
qu'il ne quitte pas plus que sa \VUgras{blouse} \textsuperscript{(47)}
décolorée, il conduit à la bergerie \VUgras{moutons} \textsuperscript{(50)},
\VUgras{brebis} \textsuperscript{(53)}, \VUgras{agneaux}
\textsuperscript{(54)}, \VUgras{chèvres} \textsuperscript{(51)} et
\VUgras{chevreaux} \textsuperscript{(52)}. Son fidèle \VUgras{chien}
\textsuperscript{(55)} Argus vient en dernier et presse de ses aboiements les
retardataires.

%%K t07-c2-11-1 p
11. --- Tout ce bruit et tout ce mouvement ne troublent pas la tranquillité de la
bonne \VUgras{vache laitière} \textsuperscript{(56)}, qui fait tinter sa
\VUgras{clochette} \textsuperscript{(57)} pendant qu'on la trait devant la porte
de l'\VUgras{étable} \textsuperscript{(58)}.

%%K t07-c2-12-1 p
12. Elle semble comprendre qu'elle doit donner son lait pour que la
\VUgras{fermière} \textsuperscript{(60)} puisse faire du \VUgras{beurre} dans la
\VUgras{baratte} \textsuperscript{(61)} ou les excellents \VUgras{fromages}
\textsuperscript{(64)} qui s'égouttent sur la planche posée sur la
\VUgras{cuve} \textsuperscript{(63)}.

%%K t07-c2-13-1 p
13. --- Toute la \VUgras{laiterie} \textsuperscript{(59)} est tenue très
proprement par le \VUgras{garçon de ferme} \textsuperscript{(65)}, qui vient de
laver les \VUgras{pots au lait} \textsuperscript{(62)} dans le grand seau qu'il
porte à la main.

%%K t07-tit-9 sub
\VUsaut{3.0mm}
\VUpk{10.2pt}{40}{La basse-cour.}
\VUsaut{3.0mm}

%%K t07-c2-14-1 p
14. --- Avec ses gros sabots de bois, il marche un peu lourdement, ce qui met en
fuite le \VUgras{dindon} \textsuperscript{(68)}, les \VUgras{canes}
\textsuperscript{(66)}, les \VUgras{canards} \textsuperscript{(67)} et les
\VUgras{oies} \textsuperscript{(69)}, qui ouvrent un large \VUgras{bec}
\textsuperscript{(70)} et crient de frayeur.

%%K t07-c2-14-2 p
Le \VUgras{coq} \textsuperscript{(71)}, plus belliqueux, dresse sa
\VUgras{crête} \textsuperscript{(72)} rouge, secoue les plumes de sa
\VUgras{queue} \textsuperscript{(73)} et fait entendre un « cocorico »
retentissant.

%%K t07-c2-15-1 p
15. --- Une \VUgras{poule mère} \textsuperscript{(74)} picore sur le
\VUgras{fumier} \textsuperscript{(81)} avec ses \VUgras{poussins}
\textsuperscript{(75)}. Le \VUgras{petit cochon} \textsuperscript{(78)} se
sauve vers sa mère, l'énorme \VUgras{truie} \textsuperscript{(77)} qu'on voit
près de la porcherie.

%%K t07-c2-16-1 p
16. --- Dans leurs \VUgras{clapiers}, les \VUgras{lapins}
\textsuperscript{(79)} mangent avidement l'\VUgras{herbe} \textsuperscript{(80)}
tendre que leur apporte la fille du fermier. Jeannine aime beaucoup ses lapins
et, chaque jour, après être allée chercher les \VUgras{œufs}
\textsuperscript{(76)} au poulailler, elle vient rendre visite à ses petits
amis.
\VUsaut{6.0mm}
\VUfilet{22mm}
